% texte.tex — f177r, vv. 1211–1231 (conversion initiale des travaux antérieurs)
\vers[n=1211]{}{}{κοὐ δύσκολ᾽ ἀλλὰ πρᾶο. ποῦ τὸ γρᾴδιο;}{et pas méchante... gentille! Où est la vieille ?}
\vers[n=1212]{}{}{οἴμ᾽ ὠς ἀπόλωλο· ποῦ τὸ γέροντ᾽ ἐντευτενί;}{Ah ! je suis perdu! Où s'est éclipsé le vieillard ?}
\vers[n=1213]{}{}{ὦ γρᾴδι᾽, ὦ γρᾶ᾽. οὐκ ἐπαινῶ γρᾴδιο.}{Ah! petite vieille, petite vieille ! Ce n'est pas bien, petite vieille.}
\vers[suite]{ΤΟ.}{Archer}{Ἀρταμουξία.}{\didasc{Appelant.} Artamouxia !}
\vers[n=1214]{}{}{διέβαλλέ μ᾽ ἡ γραῦς. ἀπότρεκ᾽ ὠς τάκιστα σύ·}{Elle m'a trompé, la vieille. \didasc{La danseuse s'esquive.} Va-t'en au plus vite, toi.}
\vers[n=1215]{}{}{ὀρτῶς δὲ συμβήνη 'στί· καταβηνῆσι γάρ.}{\didasc{Reprenant sa lance :} La lance est droite, je descends.}
\vers[n=1215]{}{}{οἴμοι, τί δρᾶσι; ποῖ τὸ γρᾴδιο;}{Ah ! que vais-je faire ? Où est la petite vieille ?}
\vers[suite]{ΤΟ.}{Archer}{Ἀρταμουξία.}{Artamouxia !}
\vers[n=1217]{ΚΟ.}{Coryphée}{τὴν γραῦν ἐρωτᾷς, ἣ 'φερεν τὰς πηκτίδας;}{Tu demandes après la vieille qui portait les pectis ?}
\vers[n=1218]{ΤΟ.}{Archer}{ναὶ ναῖκι. εἶδες αὐτό;}{Si fait, si. L'as-tu vue ?}
\vers[suite]{ΚΟ.}{Coryphée}{ταύτῃ γ᾽ οἴχεται}{C'est par là qu'elle a disparu,}
\vers[n=1219]{}{}{αὐτή τ᾽ ἐκείνη καὶ γέρων τις εἵπετο.}{elle et un vieux qui la suivait.}
\vers[n=1220]{ΤΟ.}{Archer}{κροκῶτ᾽ ἔκοντο τὴ γέροντο;}{Il avait une crocote, le vieux ?}
\vers[suite]{ΚΟ.}{Coryphée}{φήμ᾽ ἐγώ.}{Tout juste.}
\vers[n=1221]{}{}{ἔτ᾽ ἂν καταλάβοις, εἰ διώκοις ταυτῃί.}{Tu peux encore les rattraper en tirant de ce côté. \didasc{Indiquant une autre direction.}}
\vers[n=1222]{ΤΟ.}{Archer}{ὦ μιαρὸ γρᾷο· πότερα τρέξι τὴν ὀδό;}{Ô la vieille canaille! Par quel chemin prendre ?}
\vers[n=1223]{}{}{Ἀρταμουξία.}{Artamouxia ! \didasc{Il court vers l'entrée de l'orchestre.}}
\vers[n=1223]{ΚΟ.}{Coryphée}{ὀρθὴν ἄνω δίωκε. ποῖ θεῖς; οὐ πάλιν}{Monte tout droit. Où cours-tu ?}
\vers[n=1224]{}{}{τῃδὶ διώξει; ᾽ς τοὔμπαλιν τρέχεις σύ γε.}{Veux-tu bien revenir ici ? Tu cours à contresens.}
\vers[n=1225]{ΤΟ.}{Archer}{κακόδαιμον, ἀλλὰ τρέξι. Ἀρταμουξία.}{Malheureux !... Mais je me sauve. Artamouxia !}
\vers[n=1226, indent=3em, findent=1.7999999999999998em, metre={an.}]{ΚΟ.}{Coryphée}{τρέχε νυν κατ᾽ αὐτοὺς ἐς κόρακας ἐπουρίσας.}{Cours maintenant, cours aux corbeaux, et qu'un bon vent te seconde.}
\vers[n=1227, indent=3em, findent=1.7999999999999998em]{}{}{ἀλλὰ πέπαισται μετρίως ἡμῖν·}{Mais nous nous sommes assez ébattues.}
\vers[n=1228, indent=3em, findent=1.7999999999999998em]{}{}{ὥσθ᾽ ὥρα δή 'στι βαδίζειν}{Il est temps pour chacune de rentrer au logis.}
\vers[n=1229, indent=3em, findent=1.7999999999999998em]{}{}{οἴκαδ᾽ ἑκάστῃ. τὼ Θεσμοφόρω δ᾽}{Puissent, en retour, les Thesmophores}
\vers[n=1230, indent=3em, findent=1.7999999999999998em]{}{}{ἡμῖν ἀγαθὴν}{nous accorder}
\vers[n=1231, indent=3em, findent=1.7999999999999998em]{}{}{τούτων χάριν ἀνταποδοίτην.}{leur précieuse faveur !}
